\section{مقدمه}

\subsection{معرفی بیماری دیستروفی عضلانی دوشن}

دیستروفی عضلانی دوشن  (
{DMD}
)
یکی از شدیدترین و شایع‌ترین اختلالات ژنتیکی عضلانی در کودکان است که به‌صورت وابسته به کروموزوم X به ارث می‌رسد. این بیماری ناشی از جهش در ژن دیستروفین بوده و منجر به فقدان یا کاهش شدید پروتئین دیستروفین در بافت عضلانی می‌شود. دیستروفین نقش کلیدی در حفظ یکپارچگی ساختاری فیبرهای عضلانی دارد و نبود آن موجب آسیب‌پذیری سلول‌های عضلانی در برابر تنش‌های مکانیکی می‌گردد.

شروع علائم بیماری معمولاً در سنین ۳ تا ۵ سالگی مشاهده می‌شود و با ضعف تدریجی عضلات پروگزیمال، اختلال در راه رفتن، دشواری در بالا رفتن از پله‌ها و زمین خوردن‌های مکرر همراه است. با پیشرفت بیماری، ضعف عضلانی گسترش یافته و بیماران در سال‌های بعدی توانایی راه رفتن مستقل را از دست می‌دهند. روند پیش‌رونده بیماری در نهایت می‌تواند عضلات تنفسی و قلبی را نیز درگیر کند.

با توجه به ماهیت پیشرونده DMD، پایش مداوم وضعیت عملکرد حرکتی بیماران از اهمیت ویژه‌ای برخوردار است. ارزیابی دقیق تغییرات عملکردی نه‌تنها در بررسی سیر بیماری مؤثر است، بلکه در تصمیم‌گیری‌های درمانی، ارزیابی اثربخشی مداخلات دارویی و طراحی برنامه‌های توان‌بخشی نیز نقش تعیین‌کننده‌ای ایفا می‌کند.

\subsection{اهمیت ارزیابی عملکرد حرکتی و چالش‌های ارزیابی بالینی}

ارزیابی عملکرد حرکتی در بیماران مبتلا به دیستروفی عضلانی دوشن، یکی از مهم‌ترین ابزارهای پایش سیر پیشرفت بیماری محسوب می‌شود. با توجه به ماهیت پیشرونده DMD، تغییرات عملکردی بیماران به‌تدریج و در بازه‌های زمانی طولانی رخ می‌دهد؛ ازاین‌رو ثبت دقیق این تغییرات می‌تواند اطلاعات ارزشمندی در خصوص روند بیماری، پاسخ به درمان‌های دارویی و اثربخشی برنامه‌های توان‌بخشی فراهم آورد.

در حال حاضر، یکی از رایج‌ترین ابزارهای ارزیابی عملکرد حرکتی در بیماران سرپایی مبتلا به DMD، مقیاس NSAA است. این مقیاس شامل ۱۷ فعالیت عملکردی بوده و هر فعالیت بر اساس توانایی بیمار در انجام حرکت، با نمره‌ای بین صفر تا دو امتیازدهی می‌شود. نمره ۲ نشان‌دهنده انجام طبیعی و مستقل حرکت، نمره ۱ بیانگر انجام حرکت با تغییرات جبرانی یا روش اصلاح‌شده، و نمره ۰ نشان‌دهنده ناتوانی در انجام مستقل فعالیت است.

با وجود کاربرد گسترده NSAA در محیط‌های بالینی، فرآیند امتیازدهی آن مبتنی بر مشاهده مستقیم و قضاوت درمانگر است. این موضوع می‌تواند منجر به بروز تفاوت بین ارزیابان مختلف، کاهش تکرارپذیری نتایج و تأثیرپذیری از عوامل ذهنی شود. علاوه بر این، برخی تغییرات ظریف در الگوهای حرکتی ممکن است با مشاهده چشمی به‌طور کامل قابل تشخیص نباشند، درحالی‌که این تغییرات می‌توانند نشان‌دهنده مراحل اولیه افت عملکرد باشند.

از سوی دیگر، ارزیابی‌های بالینی معمولاً در فواصل زمانی مشخص و در محیط‌های درمانی انجام می‌شوند و امکان پایش پیوسته وضعیت حرکتی بیماران در شرایط روزمره را فراهم نمی‌کنند. این محدودیت‌ها ضرورت توسعه ابزارهای مکمل و داده‌محور برای ارزیابی عینی‌تر عملکرد حرکتی را برجسته می‌سازد. استفاده از حسگرهای پوشیدنی و تحلیل کمی داده‌های حرکتی می‌تواند به کاهش وابستگی به قضاوت ذهنی و افزایش دقت اندازه‌گیری کمک نماید.

در این چارچوب، بهره‌گیری از حسگرهای اینرسیایی به‌عنوان ابزاری کم‌هزینه، سبک و قابل‌حمل، امکان ثبت دقیق پارامترهای سینماتیکی نظیر شتاب خطی و سرعت زاویه‌ای را فراهم می‌کند. تحلیل این داده‌ها می‌تواند تصویری کمی از کیفیت اجرای حرکات ارائه دهد و بستری مناسب برای پیاده‌سازی روش‌های مبتنی بر یادگیری ماشین در فرآیند نمره‌دهی ایجاد کند.

\subsection{بیان مسئله}

با وجود پیشرفت‌های قابل‌توجه در حوزه درمان‌های دارویی و ژنتیکی بیماری دیستروفی عضلانی دوشن، همچنان پایش دقیق عملکرد حرکتی بیماران یکی از چالش‌های اساسی در مدیریت این بیماری محسوب می‌شود. در عمل، تصمیم‌گیری‌های درمانی، ارزیابی اثربخشی مداخلات و حتی تعیین شدت پیشرفت بیماری، تا حد زیادی بر پایه ابزارهای ارزیابی بالینی مانند مقیاس NSAA انجام می‌گیرد.

اگرچه این مقیاس به‌عنوان یک ابزار استاندارد شناخته می‌شود، اما فرآیند امتیازدهی آن مبتنی بر مشاهده مستقیم و قضاوت ارزیاب است. این وابستگی به ارزیابی انسانی می‌تواند منجر به تفاوت‌های بین‌ارزیاب، تأثیرپذیری از تجربه فردی درمانگر و محدودیت در تشخیص تغییرات ظریف حرکتی شود. علاوه بر این، بسیاری از شاخص‌های عملکردی به‌صورت کیفی تعریف شده‌اند و فاقد معیارهای عددی دقیق برای تمایز سطوح مختلف عملکرد هستند.

از سوی دیگر، پیشرفت فناوری حسگرهای پوشیدنی و توسعه سامانه‌های مبتنی بر داده، امکان اندازه‌گیری کمی پارامترهای حرکتی را فراهم ساخته است. حسگرهای اینرسیایی قادرند داده‌های دقیق شتاب و سرعت زاویه‌ای اندام‌ها را با نرخ نمونه‌برداری مناسب ثبت نمایند. با این حال، تبدیل این داده‌های خام به شاخص‌های معنادار بالینی و تطبیق آن‌ها با سیستم امتیازدهی NSAA همچنان یک مسئله باز پژوهشی محسوب می‌شود.

مسئله اصلی این پژوهش آن است که چگونه می‌توان با استفاده از داده‌های حاصل از حسگرهای اینرسیایی، شاخص‌های عملکرد موتوری بیماران مبتلا به DMD را به‌صورت کمی استخراج کرده و چارچوبی مبتنی بر یادگیری ماشین برای پیش‌بینی نمره حرکات ارائه داد؛ به‌گونه‌ای که نتایج حاصل بتوانند به‌عنوان مکملی عینی در کنار ارزیابی بالینی مورد استفاده قرار گیرند.

به بیان دیگر، این پژوهش در پی پاسخ به این پرسش است که آیا می‌توان فرآیند نمره‌دهی حرکات در مقیاس NSAA را از یک ارزیابی صرفاً کیفی و وابسته به مشاهده انسانی، به یک فرآیند داده‌محور، تکرارپذیر و مبتنی بر تحلیل کمی تبدیل نمود یا خیر.

\subsection{اهداف پژوهش}

با توجه به چالش‌های مطرح‌شده در فرآیند ارزیابی عملکرد حرکتی بیماران مبتلا به دیستروفی عضلانی دوشن، هدف کلی این پژوهش توسعه یک چارچوب داده‌محور برای ارزیابی کمی شاخص‌های عملکرد موتوری با استفاده از حسگرهای اینرسیایی و روش‌های یادگیری ماشین است.

\subsubsection*{هدف کلی}

طراحی و پیاده‌سازی سامانه‌ای مبتنی بر حسگرهای اینرسیایی به‌منظور استخراج شاخص‌های کمی حرکتی و پیش‌بینی نمره فعالیت‌های منتخب مقیاس NSAA در بیماران مبتلا به DMD.

\subsubsection*{اهداف جزئی}

\begin{itemize}
	\item طراحی و پیاده‌سازی بستر جمع‌آوری داده‌های شتاب خطی و سرعت زاویه‌ای با استفاده از حسگرهای اینرسیایی پوشیدنی.
	
	\item ثبت و ذخیره‌سازی داده‌های حرکتی بیماران در حین اجرای فعالیت‌های منتخب مقیاس NSAA.
	
	\item پیش‌پردازش داده‌های خام شامل حذف نویز، فیلترگذاری دیجیتال و نرمال‌سازی زمانی جهت آماده‌سازی برای تحلیل.
	
	\item استخراج ویژگی‌های مناسب از حوزه زمان و فرکانس به‌منظور توصیف کمی کیفیت اجرای حرکات.
	
	\item به‌کارگیری الگوریتم‌های یادگیری ماشین برای مدل‌سازی رابطه بین ویژگی‌های استخراج‌شده و نمره حرکات.
	
	\item مقایسه نتایج پیش‌بینی‌شده با ارزیابی بالینی به‌منظور بررسی دقت و قابلیت اعتماد سامانه پیشنهادی.
\end{itemize}

\subsubsection*{خروجی‌های مورد انتظار}

انتظار می‌رود سامانه توسعه‌یافته بتواند به‌عنوان ابزاری مکمل در کنار ارزیابی‌های بالینی مورد استفاده قرار گیرد و با ارائه شاخص‌های کمی، به افزایش عینیت، دقت و تکرارپذیری فرآیند نمره‌دهی کمک نماید. همچنین این چارچوب می‌تواند زمینه‌ای برای توسعه سامانه‌های پایش طولی بیماران و کاربردهای ارزیابی از راه دور در آینده فراهم سازد.

\subsection{نوآوری‌های پژوهش}

با توجه به مطالعات پیشین در حوزه تحلیل حرکت و استفاده از حسگرهای پوشیدنی، پژوهش حاضر تلاش نموده است با رویکردی یکپارچه و مسئله‌محور، گامی در جهت کمی‌سازی فرآیند ارزیابی عملکرد حرکتی بیماران مبتلا به DMD بردارد. نوآوری‌های اصلی این پژوهش را می‌توان به شرح زیر خلاصه نمود:

\begin{itemize}
	\item ارائه چارچوبی داده‌محور برای تطبیق مستقیم داده‌های حاصل از حسگرهای اینرسیایی با سیستم امتیازدهی مقیاس \lr{NSAA}، به‌گونه‌ای که نمره هر فعالیت بر اساس تحلیل کمی داده‌های حرکتی پیش‌بینی شود.
	
	\item توسعه بستر یکپارچه جمع‌آوری، پیش‌پردازش و تحلیل داده‌های حرکتی در محیط \lr{MATLAB}، شامل ماژول‌های دریافت داده، فیلترگذاری دیجیتال، قطعه‌بندی زمانی و استخراج ویژگی.
	
	\item تعریف و استخراج مجموعه‌ای از ویژگی‌های زمانی و فرکانسی که توانایی تمایز بین سطوح مختلف عملکرد حرکتی (نمرات ۰، ۱ و ۲) را داشته باشند.
	
	\item به‌کارگیری الگوریتم‌های یادگیری ماشین برای مدل‌سازی رابطه بین پارامترهای سینماتیکی و نمره بالینی، با هدف کاهش وابستگی به قضاوت ذهنی ارزیاب.
	
	\item ارزیابی عملی محدودیت‌ها و ملاحظات کاربرد سامانه در محیط بالینی، بر اساس مشاهده میدانی بیماران و تعامل با تیم تخصصی.
\end{itemize}

برخلاف بسیاری از مطالعات که صرفاً به تحلیل سیگنال‌های حرکتی یا طبقه‌بندی نوع فعالیت می‌پردازند، تمرکز این پژوهش بر کمی‌سازی کیفیت اجرای حرکت و ارتباط مستقیم آن با یک مقیاس بالینی معتبر بوده است. این رویکرد می‌تواند زمینه‌ای برای توسعه سامانه‌های ارزیابی عینی و قابل‌اتکا در کاربردهای بالینی فراهم آورد.

\subsection{ساختار پایان‌نامه}

این پایان‌نامه در قالب چند فصل سازمان‌دهی شده است تا مسیر پژوهش از معرفی مسئله تا تحلیل نتایج به‌صورت منسجم و مرحله‌به‌مرحله ارائه گردد.

در فصل دوم، مبانی نظری پژوهش شامل معرفی دقیق‌تر بیماری دیستروفی عضلانی دوشن، تشریح مقیاس \lr{NSAA}، اصول عملکرد حسگرهای اینرسیایی و مفاهیم پایه پردازش سیگنال و یادگیری ماشین بیان می‌شود. این فصل چارچوب تئوریک لازم برای درک روش پیشنهادی را فراهم می‌کند.

فصل سوم به مرور ادبیات پژوهش اختصاص دارد. در این فصل مطالعات پیشین مرتبط با تحلیل حرکت مبتنی بر حسگرهای پوشیدنی، ارزیابی کمی بیماران مبتلا به DMD و کاربرد روش‌های یادگیری ماشین در حوزه ارزیابی حرکتی بررسی و مقایسه می‌شوند.

در فصل چهارم، طراحی و پیاده‌سازی سامانه پیشنهادی شامل ساختار سخت‌افزار، نحوه نصب حسگرها، فرآیند جمع‌آوری داده، پیش‌پردازش سیگنال‌ها و استخراج ویژگی‌ها تشریح می‌گردد.

فصل پنجم به مدل‌سازی و نمره‌دهی خودکار حرکات اختصاص دارد. در این فصل الگوریتم‌های یادگیری ماشین مورد استفاده، نحوه آموزش مدل و معیارهای ارزیابی عملکرد سیستم ارائه می‌شوند.

در فصل ششم، نتایج به‌دست‌آمده تحلیل و تفسیر شده و عملکرد سامانه پیشنهادی با ارزیابی بالینی مقایسه می‌شود. در نهایت، فصل هفتم شامل جمع‌بندی، محدودیت‌های پژوهش و پیشنهادهایی برای تحقیقات آتی است.
